\begin{proposition}[B.4.13 (Proved): the second period does not lift to a finite marked Hopf modification]
Fix $0<|\alpha|,|a|<1$, and put
\[
 H=(\C^2\setminus\{0\})/
       \langle(x,v)\mapsto(\alpha x,av)\rangle,
       \qquad E=\{x=0\}\simeq\C^*/a^\Z.
\]
Let $h$ be the automorphism induced by
$(x,v)\mapsto(c x,bv)$, where $b,c\in\C^*$.
Suppose $p:Y\to H$ is a finite composition of point blow-ups,
allowing infinitely near centres, and at least one point of $E$
is an image of its exceptional locus. If $[b]\in\C^*/a^\Z$ has
infinite order, there is no holomorphic automorphism of $Y$
which agrees, on the dense unmodified core, with the marked
automorphism $h$.

The same obstruction holds for a finite collection of such marked
modifications if the proposed period permutes the collection and
acts as $h$ on each marked core.
More generally, if a compact smooth surface contains a preserved
marked elliptic curve $E$ of nonzero normal degree $d$, every
holomorphic automorphism restricting to translation by $b$ on $E$
must satisfy $d\,b=0$ in $E$.
The latter assertion also excludes a finite cycle of such preserved
marked cap axes when the cycle composite is an infinite-order
translation. These assertions apply to the direct attempt
to lift the second actual cusp period to the finite Hopf
modifications in B.4.10, for generic nonzero parameters, after every
finite base change. It concerns a holomorphic period action on
these full modifications. It does not exclude a meromorphic
correspondence, a different connected unfolding made from
noncompact pieces, or a cap present only over the central parameter.
\end{proposition}
\begin{proof}
Let $S\subset H$ be the image of the exceptional curves of $p$.
It is finite. Repeated infinitely near blow-ups change the curves
over a point but do not make $S$ infinite.
If $\widetilde h$ were a lift, the maps
\[
 p\widetilde h,\ hp:Y\longrightarrow H
\]
would agree on the dense core and hence everywhere, by the identity
principle for holomorphic maps to the separated complex surface $H$.
The image under $\widetilde h$ of a curve contracted by $p$ is
therefore again contracted by $p$. Applying the same argument to
$\widetilde h^{-1}$ proves that the total exceptional locus is
preserved. Thus $h(S)=S$.

The automorphism $h$ preserves $E$, where it is translation by
$[b]$. Consequently the nonempty finite set $S\cap E$ would be
invariant under that translation. The orbit of any one of its
points is finite, so for some $k>0$,
\[
 [b^k v]=[v],\qquad\text{equivalently } b^k\in a^\Z.
\]
This says exactly that $[b]$ is torsion, a contradiction.
In particular, no choice of multiplicities or infinitely near
centres changes the obstruction: only their nonempty finite
support on $E$ is needed.

For finitely many modifications $p_i:Y_i\to H$, a proposed deck
map has biholomorphisms
$Y_i\to Y_{\sigma(i)}$ satisfying the same core comparison.
The identity principle again gives equivariance everywhere.
The union of the finitely many exceptional image sets is
$h$-invariant; its intersection with $E$ gives the same contradiction.
Equivalently, around a cycle of length $\ell$ the map $h^\ell$
would lift to a single modification, while $[b^\ell]$ still has
infinite order.

For the more general assertion no blowdown is needed.
If an automorphism preserves $E$, its derivative in the normal
direction gives
\[
 t_b^*N_{E/S}\simeq N_{E/S}.
\]
For a line bundle $N$ of degree $d$ on an elliptic curve, the
homomorphism $b\mapsto t_b^*N\otimes N^{-1}$ is multiplication
by $-d$ under the usual identification $\operatorname{Pic}^0(E)=E$.
Indeed write $N=\OO_E(\sum_j n_j[p_j])$ with
$\sum_j n_j=d$; the Abel--Jacobi class of its translated divisor
minus the original one is $-d b$.
Thus preservation of $N$ forces $d b=0$.
Taking a power proves the finite-cycle version, since its composite
preserves one axis and its nonzero-degree normal line.
In particular a nonzero-degree parabolic axis cannot carry the
marked generic period translation even if no finite Hopf
modification is specified. The original finite-support argument
also covers modifications where that normal-degree test alone
would give no information.

We check the actual cusp rather than assume that torsion fails.
The periods of A.1.2 have the form
\[
 g_1:(Z,W)\longmapsto(q c_1 Z,c_2W),\qquad
 g_2:(Z,W)\longmapsto(c_3Z,r^{12}q c_4 W),
\]
where the $c_i$ are nonvanishing holomorphic units near the cusp.
Choose $W$ or $W^{-1}$ so that
$a=c_2^{\varepsilon}$ has modulus less than one,
$\varepsilon\in\{1,-1\}$.
This is possible at the actual marked cusp because its limiting
unit has modulus different from one, as computed in B.4.8.
For $q\ne0$ small, the partial $g_1$ quotient of $(\C^*)^2$
is the dense core of the diagonal Hopf surface with multipliers
$(q c_1,a)$. The second period induces the above $h$, with
\[
 b(q,r)=(r^{12}q c_4(q,r))^{\varepsilon}.
\]
Fix a sufficiently small nonzero $q$. For each $k>0$ and
$j\in\Z$, the equation $b(q,r)^k=a(q)^j$ has a discrete
solution set in the punctured $r$-disc: its left-hand side has
a zero or pole of nonzero order at $r=0$, while the right-hand
side is a nonzero constant in $r$.
Hence outside a countable union of discrete sets, $[b(q,r)]$
has infinite order. Under a finite base change in $r$ the
nonzero order is multiplied by its ramification index and
the same conclusion holds. A finite base change in $q$ does
not affect this fixed-nonzero-$q$ argument.

To compare with B.4.10, write $q c_1=\kappa(s)s^n$ after
the chosen finite cover in the cusp variable. Its generic Kato
fibre is precisely a finite Hopf modification with nonempty
axis centre $[-s]$. A holomorphic extension of the actual second
period to that whole fibre would therefore contradict the result
just proved. Allowing finitely many translated copies does not
remove the contradiction, by the finite-collection assertion.
\end{proof}

\begin{remark}[What an infinite unfolding must additionally prove]
There is no corresponding obstruction to merely indexing translated
modifications by all integers and letting the deck map shift the
index. That construction alone has the wrong core. Indeed suppose
$p_j:Y_j\to H$ and
$\widetilde h_j:Y_j\to Y_{j+1}$ satisfy
$p_{j+1}\widetilde h_j=hp_j$.
On the disjoint union of their unmodified cores, the quotient by
the shift is identified with the single core $H^\circ$ by
\[
 [(j,y)]\longmapsto h^{-j}p_j(y).
\]
It is not the required quotient $H^\circ/\langle h\rangle$.
Thus a useful unfolding must have additional connected overlaps
and must recover the original second-period identifications there.

Nor can one simply glue several full compact Hopf modifications
along the same dense core and assume a Hausdorff space results.
If two embedded compact spaces in a Hausdorff space share an open
subset dense in both, their images, being closed, both equal its
closure. They coincide. A construction by genuinely distinct
compact copies with only that common core consequently cannot be
such a Hausdorff gluing.
These observations leave open connected unfoldings of noncompact
pieces with a new boundary geometry. Their analytic transition
maps, proper discontinuity, correct punctured quotient and compact
central quotient would all have to be proved; none follows from
the auxiliary compact family in B.4.10.
\end{remark}
